\chapter{Mission and instrument context}
\label{chapter:context}

The PLASMO design report holds a first set of formal requirements. All values are current best estimates and several carry TBC flags, so the deployment design has to remain adaptable. The values most relevant to the M2 structure are given in \Cref{tab:requirements}.

\begin{table}[H]
    \centering
    \caption{Requirements relevant to the deployable M2 structure, from the PLASMO design report, July 2026 snapshot \cite{bouwmeester_plasmo_2026}. Values are current best estimates and carry TBC flags.}
    \label{tab:requirements}
    \small
    \begin{tabularx}{\textwidth}{lXX}
        \toprule
        Item & Current value & Consequence for the M2 structure \\
        \midrule
        
        Ground sampling distance & $\leq$ \SI{10}{\meter} & Drives the aperture, and through it the M2 size and standoff \\
        \addlinespace
        
        Swath & $\geq$ \SI{10}{\kilo\meter} & Field of view; couples to the optical layout \\
        
        Spectral bands & Seven VIS-SWIR bands in four polarisations & Acousto-optic filtering; instrument volume \\
        \addlinespace
        
        SNR & $\geq$ 100 & The system is SNR-limited, which pushes the aperture up \\
        \addlinespace
        
        Orbit & 200 to \SI{350}{\kilo\meter}, \SI{300}{\kilo\meter} baseline, SSO & VLEO thermal cycling and atomic oxygen \\
        \addlinespace
        
        Aperture & At least 20 to \SI{30}{\centi\meter}; radiometry recommends \SI{30}{\centi\meter} & Design at the larger end; the standoff grows with it \\
        \addlinespace
        
        M1 to M2 distance & 1.5 to 2 times the M2 diameter, or in absolute terms ideally around \SI{30}{\centi\meter} and not more than \SI{60}{\centi\meter} & The main open variable; see \Cref{sec:families-assessed} and \Cref{sec:questions} \\
        \addlinespace
        
        Optical configuration & Three-mirror folded design, the third mirror turning the light through 90 degrees into the instrument box. On-axis preferred, off-axis not excluded & M2 stays the only deployed optic; the preference narrows sub-question 2.2 without closing it \\
        \addlinespace
        
        Deployment accuracy & No requirement defined yet & The accuracy budget is open and is expected to be set in the design report first \\
        \bottomrule
    \end{tabularx}
\end{table}

Three consequences follow from these requirements.

First, currently no deployment-accuracy requirement is defined, so the thesis derives its own target rather than inheriting one, which is why the accuracy budget is a deliverable in its own right (\Cref{sec:phase2}). Its inputs and assumptions are stated alongside the outputs of the analyses, so that the budget can be updated and fed back to the design report as the instrument design settles.

Second, different architectures could be considered, depending on the separation between M1 and M2. There will be M1-M2 separation distances where a specific architecture is preferable over another. For example, folding architectures, whose arms must lie against the bus, are preferable over extending architectures up to a separation of roughly the length of the bus size. Since the bus size is probably determined by the primary mirror diameter, beyond this extending architectures could take over. For a \SI{30}{\centi\meter} primary mirror, this means folding architectures work well up to about \SI{30}{\centi\meter} separation, but are no longer suitable beyond that. The architectures considered in this paper are set out in \Cref{sec:families-assessed}.

Third, the upper bound on the separation is a working guideline rather than a derived limit. The optics push the separation up, since a high F number is needed to stay within the AOTF acceptance angles, and the design report pushes back on physical length for two platform reasons: the projected area of the baffle in the velocity direction, which drives drag in VLEO, and the thermo-mechanical stability of a deployable telescope \cite{bouwmeester_plasmo_2026}. The \SI{30}{\centi\meter} ideal is the separation that earlier drag calculations already assumed, and the \SI{60}{\centi\meter} cm ceiling is a two to one ratio on a \SI{30}{\centi\meter} primary, set as a target for the optical design. Both scale with the aperture, which is itself only bounded to \SI{20}{\centi\meter} to \SI{40}{\centi\meter}, and both may move once an optical prescription and a drag budget exist. The separation is therefore swept up to 60 cm and a little beyond, so that the models also show what the architectures cost if the optical design returns a longer instrument; 60 cm is carried as a reference ceiling, not as a hard constraint. Drag itself enters at system level, as a penalty on instrument length, rather than as a load case on the deployed structure, because the baffle shields M2 and its support.